\clearpage
\fancyhead[L]{\small ZWEITER COMET VOM JAHR 1805.}
\hypertarget{native-280-lower}{}
\begin{center}
ZWEITER COMET VOM JAHR 1805.

\rule{0.6\linewidth}{0.4pt}

Monatliche Correspondenz. Band XIV. S. 181–186. 1806 August.

\rule{0.6\linewidth}{0.4pt}
\end{center}

\begin{center}
\textsc{Gauss an von Zach.}
\end{center}

\begin{flushright}
Braunschweig, den 8. Juli 1806.
\end{flushright}

\noindent --- --- --- Über den Cometen von 1805 habe ich nach dem, was ich Ihnen bereits geschrieben habe, nichts weiter gearbeitet, als die Vergleichung der mir gütigst mitgetheilten \textsc{Bouvard}'schen Beobachtungen mit meinen Elementen, sowol den parabolischen (M.\,C. Juli-Heft 1806. S. 77) [d.\,B. S. 271] als den elliptischen (ebendas. S. 79) [d.\,B. S. 272]. Diese Vergleichung steht so: [Differenz ist R---B.]

\begin{center}\begingroup
\fontsize{8}{10}\selectfont
\setlength{\tabcolsep}{1pt}\renewcommand{\arraystretch}{1.12}
\begin{tabular}{l||c|c|c|c||c|c|c|c}
\shortstack{Zeit d. Beobacht.} & \multicolumn{2}{c|}{\shortstack{In d. Parabel\\berechnete}} & \multicolumn{2}{c||}{\shortstack{Differenz\\in der}} & \multicolumn{2}{c|}{\shortstack{In d. Ellipse\\berechnete}} & \multicolumn{2}{c}{\shortstack{Differenz\\in der}} \\
Paris 1805. & Länge & Breite & Länge & Breite & Länge & Breite & Länge & Breite \\
\hline
Nov. 16.\,45203 & $29^\circ\,21'\,30''$ & $30^\circ\,2'\,32''$ & $-2'\,7''$ & $-8''$ & $29^\circ\,22'\,48''$ & $30^\circ\,1'\,18''$ & $-0'\,49''$ & $-1'\,22''$ \\
17.\,36386 & $28\,51\,5$ & $29\,49\,40$ & $-0\,18$ & $+48$ & $28\,51\,55$ & $29\,48\,47$ & $+0\,32$ & $-0\,4$ \\
18.\,39927 & $28\,14\,28$ & $29\,32\,42$ & $+0\,10$ & $-2$ & $28\,14\,43$ & $29\,32\,8$ & $+0\,25$ & $-0\,36$ \\
23.\,32241 & $24\,41\,33$ & $27\,25\,5$ & $+0\,29$ & $-30$ & $24\,39\,54$ & $27\,25\,31$ & $-1\,10$ & $-0\,4$ \\
30.\,51095 & $15\,43\,40$ & $19\,25\,32$ & $+4\,0$ & $+4$ & $15\,42\,3$ & $19\,26\,26$ & $+2\,23$ & $+0\,58$ \\
Dec. 5.\,29581 & $2\,7\,19$ & $3\,20\,5$ & $+0\,8$ & $-40$ & $2\,7\,20$ & $3\,21\,26$ & $+0\,9$ & $+0\,40$ \\
\end{tabular}\endgroup\end{center}

\par\indent Diese Vergleichung zeigt, dass die Ellipse mit diesen Beobachtungen nicht ganz so gut stimmt, wie mit den Marseiller Beobachtungen; allein es scheint zugleich zu erhellen, 1) dass letztere Beobachtungen besser sind, und 2) dass auch bei erstern die Parabel grössere Differenzen gibt, als die Ellipse. Doch muss ich hiebei noch bemerken, dass die Differenzen, die hier bei der Ellipse angesetzt sind, eigentlich nicht die wahren sind, sondern noch ansehnlicher Veränderungen bedürfen. Die \emph{berechneten} Längen und Breiten sind hier nemlich nicht mit Parallaxe und Aberration behaftet, und diejenigen, womit ich sie verglichen habe, sind die beobachteten, nachdem sie von Parallaxe und Aberration befreit sind (so wie Sie mir solche mitgetheilt haben). Diese Reduction, die schon gemacht war, hängt aber von den Abständen des Cometen von der Erde ab, und zwar steht die Reduction, die von der Parallaxe abhängt, in umgekehrtem, hingegen die, welche von der Aberration abhängt, in geradem Verhältnisse des Abstandes. Nun aber fallen die Abstände in der Parabel und Ellipse sehr verschieden aus; so ist z.\,B.

\begin{center}\begingroup
\fontsize{9}{11}\selectfont
\setlength{\tabcolsep}{2pt}\renewcommand{\arraystretch}{1.12}
\begin{tabular}{rlll}
Abstand in der Parabel: & Nov. 16 $=0.2489$, & Nov. 23 $=0.1761$, & Dec. 5 $=0.0638$ \\
in der Ellipse: & $=0.1794$ & $=0.1269$ & $=0.0460$ \\
\end{tabular}\endgroup\end{center}
\SourceLeaf{281}%
\par\indent Die Reduction der Parallaxe, die hier vorzüglich beträchtlich ist, fällt also in der Ellipse ungefähr in dem Verhältnisse von 7:5 grösser aus, als in der Parabel, hingegen die wegen der Aberration in demselben Verhältnisse kleiner. Da nun in den Beobachtungen diese beiden Reductionen nicht einzeln angegeben sind, sondern sich aus den Angaben nur ihre Summe abnehmen lässt, und zwar auch nur bei den vier letzten, so müsste man die Reductionen für die Ellipse von neuem berechnen, welche Mühe ich mir nicht gegeben habe. Die Reductionen in der Parabel hingegen werden eben so ausfallen, wie \textsc{Legendre} sie angenommen hat, da meine parabolischen Elemente so wenig von den seinigen abweichen. So viel kann man indessen durch blosse Schätzung schon übersehen, dass die Beobachtung vom 5. December, die einzige, welche eine Meridian-Beobachtung zu sein scheint, durch die genauere Reduction noch viel besser mit den elliptischen Elementen harmoniren würde. Sollte ich noch \textsc{Bouvard}'s Original-Rectascensionen und Declinationen zu Gesichte bekommen, so werde ich diese noch auf dieselbe Art, wie ich es bei den \textsc{Thulis}'schen gethan habe, mit meinen Elementen vergleichen. Vor der Hand habe ich diese Rechnungen bei Seite gelegt, da sich doch nichts entschiedeneres über den Cometen wird sagen lassen, so lange nicht Beobachtungen bekannt werden, die den Zeitraum noch erweitern.

\par\indent \textsc{Legendre}'s Werk, das Sie bei dieser Gelegenheit erwähnen, habe ich noch nicht gesehen. Ich hatte mit Fleiss mir deswegen keine Mühe gegeben, um bei der Arbeit an meiner Methode ganz in der Kette meiner eigenen Ideen zu bleiben. Durch ein paar Worte, die \textsc{de la Lande} in der letztern \emph{Histoire de l'Astronomie}, 1805, fallen lässt, \emph{méthode des moindres quarrés}, gerathe ich auf die Vermuthung, dass ein Grundsatz, dessen ich mich schon seit zwölf Jahren bei mancherlei Rechnungen bedient habe, und den ich auch in meinem Werke mit gebrauchen werde, ob er wol zu meiner Methode eben nicht wesentlich gehört, --- dass dieser Grundsatz auch von \textsc{Legendre} benutzt ist. Übrigens können eben nicht viele Berührungspunkte vorkommen, da unsere Arbeiten sehr verschiedene Gegenstände haben, denn die Bestimmung einer Cometenbahn, wo man eine Parabel voraussetzt, und die Bestimmung einer Planetenbahn, wo man von aller Hypothese abstrahiren soll, erfordern eine ganz verschiedene Behandlung. Letzterm Probleme ist aber mein Werk \emph{eigentlich} gewidmet, obgleich vieles auch über das erstere mit vorkommen wird, was mir eigenthümlich ist.

\par\indent Es ist mir übrigens überaus lieb, dass ich nicht schon 1802 meine Methode, wie ich die Ceres- und Pallas-Bahn berechnet hatte, bekannt gemacht habe, so viele Aufforderungen auch deshalb an mich gelangten. Denn seitdem habe ich noch immer an der Vervollkommnung der Methode selbst gearbeitet, besonders in dem vorigen Winter, und ihre jetzige Gestalt sieht ihrer ersten fast gar nicht mehr ähnlich. Um hievon eine Probe zu geben, will ich nur \emph{eines} Umstandes erwähnen. Da das Problem so sehr verwickelt ist, so ist es der Natur der Sache nach nicht anders möglich, als dass bei der allerersten Annäherung einige Voraussetzungen gemacht werden müssen, die nur näherungsweise richtig sind (wie z. B. bei Dr. \textsc{Olbers}' Methode die ist, dass die Chorden bei der Erde und dem Cometen durch die mittleren radii vectores im Verhältnisse der Zwischenzeiten geschnitten werden). Voraussetzungen von dieser Art liegen also auch nach meiner Methode bei der ersten Annäherung zum Grunde, und zwar solche, die desto weniger von der Wahrheit abweichen, je kleiner die Zwischenzeiten sind. Man darf also die Beobachtungen, auf die man die erste Annäherung gründet, nicht gar zu weit von einander entfernt annehmen, weil man sonst vermöge der näherungsweise wahren Voraussetzung bei der ersten Rechnung gar zu weit von der Wahrheit zurückbleiben, und daher zu viele und beschwerliche Wiederholungen der Verbesserungs-Methoden machen müsste. Doch konnte ich, wie meine Methode 1802 war, bei der Pallas sogleich Beobachtungen anwenden, die 27 Tage auseinander waren; viel weiter hätte ich indess doch nicht gehen mögen. Dagegen ist jetzt meine Methode so beschaffen, dass ich neulich, als ich die mir von Ihnen gütigst mitgetheilten Beobachtungen \textsc{Oriani}'s von 1805 zu einem für mein Werk
\SourceLeaf{282}%
bestimmten \emph{Exempel} benutzen wollte, und also dieselben so behandeln musste, als wenn ich von der
Pallas-Bahn noch gar nichts wüsste, sogleich und zwar mit dem allerglücklichsten Erfolge die äussersten
71 Tage von einander entfernten Beobachtungen zum Grunde legen konnte, und es leidet gar keinen
Zweifel, dass ich darin noch beträchtlich weiter hätte gehen können.

\Needspace{18\baselineskip}
\fancyhead[L]{\small GAUSS AN VON ZACH.}
\begin{center}
GAUSS AN VON ZACH.
\end{center}

\begin{center}
Monatliche Correspondenz. Band XIV. S. 187–192. 1806 August.
\end{center}

\begin{flushright}
Braunschweig, den 8. Juli 1806.
\end{flushright}

\par\indent --- --- --- Mit Hülfe der mir gütigst überschickten Mailänder Beobachtungen des Planeten
\emph{Pallas} (M. C. Juli-Heft 1806. S. 90) habe ich neue Elemente (IX.) dieser Planeten-Bahn berechnet, die
sich an die \textsc{Oriani}'schen Beobachtungen möglichst genau anschliessen, und aus diesen die Opposition
hergeleitet, was mit den zu stark abweichenden VIII. Elementen nicht so gut hätte geschehen können;
nemlich:

\begin{center}
1805. Novbr. 29.\ \ $11^{\mathrm{u}}\,18^{\mathrm{m}}\,17^{\mathrm{s}}$ M. Z. in Seeberg.

Länge:\quad $67^\circ\,20'\,42''9$

südl. geoc. Breite:\quad $54\;\;30\;\;54.9$
\end{center}

\par\indent Zum Behuf einer Ephemeride für die nächste Erscheinung der \emph{Pallas} habe ich nun auf die drei
Oppositionen von 1803, 1804, 1805 und die Beobachtungen von 1802 meine neuen Elemente gegründet.
Es zeigte sich hiebei, dass diese verschiedenen Örter sich nicht mehr ganz genau durch eine reine Ellipse darstellen liessen, was ohne Zweifel Folge der noch vernachlässigten Störungen ist, doch äussert
sich diese Discordanz immer noch als eine Kleinigkeit, die kaum eine Minute übersteigen mag. Namentlich
folgt aus den Beobachtungen von 1803 die Neigung der Bahn $44''$ grösser, und aus denen von 1805 um
so viel kleiner, als das bei den folgenden Elementen angesetzte Mittel, und der Knoten erscheint nach
den Beobachtungen von 1802 etwa $45''$ weiter zurück, und nach denen von 1804 nur eben so viel weiter
vorwärts, als in dem angesetzten Mittel. Hier sind die neuen Elemente (IX.), wonach die folgende Ephemeride berechnet worden ist:

\Needspace{19\baselineskip}
\begin{center}\begingroup
\fontsize{9}{11}\selectfont
\setlength{\tabcolsep}{2pt}\renewcommand{\arraystretch}{1.12}
\begin{tabular}{lr}
\multicolumn{2}{c}{Epoche im Meridian von Seeberg} \\
1802. & $143^\circ\,24'\,1''1$ \\
1803. & $221\;32\;56.4$ \\
1804. & $299\;54\;42.5$ \\
1805. & $18\;3\;37.8$ \\
1806. & $96\;12\;33.2$ \\
1807. & $174\;21\;18.5$ \\
\end{tabular}\endgroup\end{center}
\begin{center}\begingroup
\fontsize{9}{11}\selectfont
\setlength{\tabcolsep}{2pt}\renewcommand{\arraystretch}{1.12}
\begin{tabular}{lr}
tägliche mittlere tropische Bewegung .\ .\ . & $770''7816$ \\
Länge der Sonnenferne 1805 .\ .\ .\ . & $301^\circ\,11'\,31''1$ \\
Länge des aufsteigenden Knotens 1805 . & $172\;30\;5.8$ \\
Excentricität .\ .\ .\ .\ .\ .\ .\ .\ .\ .\ . & $0.2453840$ \\
Logarithm der mittlern Entfernung .\ .\ . & $0.4421021$ \\
Neigung der Bahn .\ .\ .\ .\ .\ .\ .\ . & $34^\circ\,37'\,8''1$ \\
\end{tabular}\endgroup\end{center}

\par\indent Die nächste Opposition würde nach diesen Elementen 1807 den 4. Mai um 21 Uhr in $223^\circ\,53'$ Länge,
und $42^\circ\,18'$ nördl. geoc. Breite eintreten. Die Zeit wird lehren, in wie fern diese Rechnung zutreffen
wird. Meine Ephemeride habe ich sogleich an Hrn. Prof. \textsc{Harding} geschickt, er wird darnach eine
Karte zeichnen, und ich hoffe, Sie werden uns bald durch die M. C. mit derselben beschenken.
\Needspace{15\baselineskip}
\SourceLeaf{283}%
\par\indent Meine nächste Arbeit wird nun die \emph{Juno} betreffen; viel wird an meinen V. Elementen nicht zu
verbessern sein. Ich habe auch diese Rechnung bereits angefangen. Hier ist vorläufig meine Vergleichung der V. Elemente mit den drei letztern von Hrn. \textsc{Bessel} im Mai angestellten Beobachtungen:

\begin{center}\begingroup
\fontsize{9}{11}\selectfont
\setlength{\tabcolsep}{2pt}\renewcommand{\arraystretch}{1.12}
\begin{tabular}{l|l|l|l|l}
& \multicolumn{2}{c|}{Berechnete} & \multicolumn{2}{c}{Unterschied in der} \\
\shortstack{Mittlere Zeit\\in Lilienthal} & \shortstack{gerade\\Aufsteigung} & \shortstack{nördliche\\Abweichung} & \shortstack{geraden\\Aufsteigung} & Abweichung \\
\hline
1806. Mai 17.\ $11^{\mathrm{u}}\,47^{\mathrm{m}}\,59^{\mathrm{s}}$ & $165^\circ\;4'\;11''$ & $9^\circ\;30'\;36''$ & $+0'\,51''$ & $+38''$ \\
19.\ $11\;\;1\;\;38$ & $165\;\;16\;\;34$ & $9\;\;29\;\;23$ & $+0\;\;33$ & $-4$ \\
23.\ $11\;\;31\;\;\;2$ & $165\;\;44\;\;57$ & $9\;\;25\;\;13$ & $+1\;\;11$ & $0$ \\
\cline{3-5}
& & Mittel & $+52''$ & $+11''$ \\
\end{tabular}\endgroup\end{center}

\par\indent Sobald ich die Elemente der \emph{Juno} verbessert habe, werde ich sogleich auch für diese eine Ephemeride berechnen und sie Ihnen und Hrn. Prof. \textsc{Harding} mittheilen. Zuletzt werde ich dieselbe Arbeit
noch für die \emph{Ceres} vornehmen, da ich hoffe, dass bis dahin noch mehrere Beobachtungen um die Zeit
der Opposition bekannt werden, als bis jetzt da sind. [Der geocentrische Lauf der Pallas nach den IX.
Elementen ist vom 7. Dec. 1806 bis 21. Sept. 1807 mit drei Tagen Zwischenräumen für 12 U. Mitternacht
in Seeberg angegeben, die gerade Aufsteigung und Abweichung in Minuten der Abstand von der Erde
in drei, die Lichtstärke in vier Decimalstellen.]
